% !TeX root = ../papers/nc-p3.tex
\section{滑らかな七次点曲線の全係数と五次元の同型類}
\label{sec:septic-bielliptic-forms}
\begin{theorem}\label{thm:septic-bielliptic-forms}
前節の滑らかな七次枝の有限条件は次のように解ける。$R=r^2,S=s^2,H=h^2+\lambda rs$ と書き、
射影四空間の座標 $r,s,h,u,v$ における三式を
\[
 ru=F=f_S S+f_H H,\qquad sv=G=g_R R+g_H H,\qquad
 v^2-u^2=4D=4(d_R R+d_S S+d_H H)
\]
とする。 $\lambda$ を含む八係数の許容範囲は三式が滑らかな完全交叉曲線を定める開集合全体である。
二点は $r=s=h=0,\ v=\pm u$ であり、前節の順序付き二射影から代数を作る。
同型の全同一視は
\[
 (r,s,h,u,v)\longmapsto (a r,b s,c h,e u,e v),\qquad
 (r,s,h,u,v)\longmapsto(s,r,h,v,u)
\]
から誘導される係数変更で尽くされる。後者は $D$ の符号も変える。
座標の共通倍率を除く三次元トーラスの重み作用なので、零係数を含め台とその重みの
整数核単項式を比較し、最後に交換を試せばよい。
この開集合は非空で同型類の次元は5である。最小周期1の必要十分条件は $\lambda$ だけを $-\lambda$ にした係数が同じ軌道に入ることであり、一般点の最小周期は2である。
全曲線が $h$ の符号反転で滑らかな楕円曲線の二重被覆になる。
\end{theorem}

\begin{proof}[標準二次式の制約]
前節の順序付き曲線から始める。三つの共通文字の空間を $J$ とし、
第一の補文字を $x$、第二を $y$ として標準空間を $J\oplus kx\oplus ky$ と書く。
標準像の二次式の基底を
\[
 Q_j=xB_j+C_j,\qquad B_j=a_j y+b_j,\qquad C_j=y c_j+d_j
 \quad(b_j,c_j\in J,\ d_j\in\operatorname{Sym}^2 J)
\]
とする。両補文字の軸が二つの指定点である。
$B_j$ は $x$ 軸での勾配なので独立である（前節の余法空間への単射）。
全 $a_j$ が零なら二辺の三文字直交点の対が三つの可換関係と全 $Q_j$ を消す。
その対は相異なる二つの曲線点に由来するのでグラフ性に反する。
従って $a_1=1,a_2=a_3=0$ とできる。$b_2,b_3$ と $c_2,c_3$ はそれぞれ独立である。
曲線の標準像はこの三式の完全交叉である。実際 $J$ 座標が非零の開集合では
二つの射影像での $J$ 座標を同一に規格化でき、三式の解がそのまま点グラフに入る。
外に残る直線上の解は二指定点だけであり、完全交叉の次数8と標準像の次数が一致する。

第一三次重なりを外二文字の反転で分解する。反対称部分では、三つの $Q_j$ に掛かる
外文字 $f_j$ は $\sum Q_j f_j=0$ を満たす。完全交叉には一次係数のsyzygyがないので
全て零である。残りの重なりは三つの可換関係だけを両端に持ち
$\bigwedge^3J$ で尽きる。従って対称部分は三次元である。
その元を最初の二文字の六関係の右倍で表し、$Q_j$ の係数を $f_j=A_jx+w_j$、
$A_j\in k,w_j\in J$ とする。対称部分から $(f_j)$ への射は単射である。
核なら上の $\bigwedge^3J$ に入り、これは反対称だからである。

対称性を中央文字 $y$ で縮約し、外側の $x$ と $J$ の成分を比較すると
$\sum a_j w_j=\sum A_j c_j$、また $\sum c_j\wedge w_j=0$ となる。
中央を $J^*$ で縮約した同じ比較から、$J\otimes J$ の対称商では
$\sum b_j w_j=\sum A_j d_j$ となる。可換関係の $x$ 倍はこの対称商で消える。
従って $\sum B_j w_j=\sum A_j C_j$ であり、
$\sum A_j Q_j$ は $(B_1,B_2,B_3)$ に入る。この三一次式の共通直線は $x$ 軸の点での接線である。
接線は完全交叉曲線に含まれず、従って $(A_j)$ の像は高々二次元である。
三次元の対称部分には $A_j=0$ の非零元がある。上の式から $w_1=0$、
$(w_2,w_3)=t(b_3,-b_2)$、$t\ne0$ となる。従って
$c_2\wedge b_3-c_3\wedge b_2=0$。$b_2,b_3$ を基底に延ばして係数を比べると、
$c_2,c_3$ は同じ二平面にあり、$b_j\mapsto c_j$ の跡は零である。
この可逆写像は対角化でき、二固有値は互いに逆符号である。

二固有文字を $r,s$ に取り、$y$ を再尺度化する。第一式の $x,y$ の一次交差項を
それぞれ $y,x$ の $J$ による持上げ変更で吸収すれば、三標準式は
\[
 xy=D,\qquad xr-r y=F,\qquad xs+s y=G
\]
となる。ここではまだ $D,F,G$ は $r,s,h$ の任意の二次式である。
非可換な二窓ではこの順の式と逆順の式を交互に置き、$r,s,h$ は交換する。
次数3だけの書換えを考える。$r<s<h<x$（次位置では $x$ を $y$ に替える）とし、
各可換逆順対と補文字に続く $r,s$, 補文字対を解消する。
次数辞書式順序は各書換えで下がる。先頭単項式の三文字重複は
$hsr,xsr,xyr,xys,xyx$ の五個である。最初の二個の差は零、
残る三個の差は非零の定数倍を除き
$xF-Fx,xG-Gx,xD-Dx$ の書換え結果である。これは $yr=rx-F,\ ys=-sx+G$ を代入すれば従う。
二次先頭関係の左右倍の行消去では各重複に一つの差が残るから、
三次重なりの次元は5からこれらの差の階数を引いた数である。

$T\in\operatorname{Sym}^2J$ に対する $xT-Tx$ の書換えが零となる核は
$\langle r^2,s^2\rangle$ である。二つの平方は直接交換する。他の四単項式
$rs,rh,sh,h^2$ の結果にはそれぞれ $-2rsx,\ ryh,\ -syh,\ xh^2$ が現れ、
その各語は他の三つの結果には現れない。従って $D,F,G$ のこの核を法とした像は一本の直線を張る。
共通の代表二次式で $h^2$ の係数が零なら、標準空間の $h$ 軸の点が曲線上にあり、
三式のそこでの勾配は高々一方向となって滑らかさに反する。
平方完成で $h$ の $r,s$ による持上げを変更し、代表の $rh,sh$ 項を消す。
残る $rs$ 項を保持して共通代表を $H=h^2+\lambda rs$ とする。
$u=x-y, v=x+y$ に移り、$u$ に $r$ の倍、$v$ に $s$ の倍を加えて
$F$ の $R$ 係数、$G$ の $S$ 係数を零にする。この操作で $D$ に加わる項も $R,S,H$ の線型結合である。
\end{proof}

\begin{proof}[滑らかな全係数での正則性]
表示された滑らかな完全交叉は連結で、随伴で超平面束が標準束、種数が5である。
長さ3の部分スキームを含む直線があれば三二次式はその直線で恒等的に消え、整曲線性に反する。
よって二指定点からの射影は接方向を含めて閉埋込みになる。
三つの可換式と三つの混合式が定める点対応もこの曲線に等しい。
両側の $J$ 座標が非零の部分は既に見た。片側だけ零のところではもう一方の $J$ 座標の
一つを1にし、零となる側の補文字を1にする。可換式により小さい側の三座標は
他方の三座標の $t$ 倍である。三混合式は補文字軸の点での標準三式の
厳密変換（代入して $t$ を一つ除いた式）で、$t=0$ での式は独立な三勾配である。
従って例外ファイバーはその接線方向の一点で、その三つのアフィン方向変数に関する
Jacobianは可逆となる。ここでも滑らかな曲線で余分な構造はない。
両側で $J=0$ は $xy$ の式に反する。

以下 $x_i$ は偶数位置の $x$ と奇数位置の $y$ を表す。関係は
$x_i r=r x_{i+1}+(-1)^i F,\ x_i s=-s x_{i+1}+G,\ x_i x_{i+1}=D$ と書ける。
二移動の意味で $R,S$ は中心的である。$r,s,h$ は交換し、$x_i r^2, x_i s^2$
を二度書き換えると補正が相殺する。$x_i x_{i+1}r, x_i x_{i+1}s, x_i x_{i+1} x_{i+2}$
の両括弧付けは $F,G,D$ のそれぞれと $x$ の交換を与える。
滑らかさから $d_H,f_H,g_H$ の少なくとも一つは非零で、$H$ も中心的である。
抽象的な次数2の中央多項式環 $B=k[R,S,H]$ 上で、$r^2=R,s^2=S,h^2=H-\lambda rs$
と上の三つの $x_i$ の式を課すと元代数に戻る。$B$ 上の基底は
$r^a s^b$ ($0\le a,b\le1$) に $h,x$ の交互語を続けたものである。
実際、$r,s$ を左に移し、平方と補文字対を消す。$r,s,h$ の内部では全て交換し、
平方の値もそこで交換する。補文字を含む重複は $xrr,xss,xsr,xxr,xxs,xxx$ で、
最初の二つは $F,G$ と $r,s$ の交換、三番目は $r,s$ の交換と
$F$ の交代符号、最後の三つは $F,G,D$ の中心性により一致する。
次数辞書式で減少するためこれで基底が従い、Hilbert級数は
$(1+z)^2(1+z)/((1-z)(1-z^2)^3)=(1-z)^{-4}$ である。

二次双対を $T$ とする。先の三次計算は全四位置で $\dim T_3=4$ を与える。
ここで下付きは各始点からの長さとする。また $\dim T_4\ge1$ である。
実際、単純の極小解消の初めは線型で生成・関係数が4,6、
三次の最小syzygyは4個である。四次のEuler式は
$\dim\operatorname{Tor}_{4,4}-\dim\operatorname{Tor}_{3,4}=1$ を与え、
対角Torの双対は四次の重なりである。

$T$ の積を具体的に調べるため、元代数の奇数位置だけ $s\mapsto-s,\ y\mapsto-X$ とし、
偶数では $x=X$ とする。双対の文字を小文字 $r,s,h,z$（最後が $X$ の双対）で再び書く。
以下の関係は六式との係数対合で読める。元の $J$ の三可換式は
$r,s$ と $s,h$ について対称、$r,h$ について交代となる。残る三式の
$X$ を含む項は $X^2,\ Xr+rX,\ Xs+sX$ であり、その右辺は
$r^2,s^2,h^2\pm\lambda rs$ の線型結合である。最後の $rs$ は $r,s$ の対称関係を法として
交代表示に取れ、符号だけが位置で交互となる。従って双対で
\begin{itemize}
\item $zh=hz=0$、$h$ は $r$ と符号を替えて、$s$ とはそのまま交換する。
$z$ は $r,s$ と交換する。$rs-sr$ は位置別スカラー倍の $h^2$ である。
\item $r^2,s^2,h^2$ はそれぞれ $z\langle z,r,s\rangle$ で表せる。
これら三つの表示は位置に依存しない。従って $T/(h)$ は
$k[r,s,z]$ の三つの二次式による商の反復である。
\end{itemize}
この三式に共通の射影零点があれば、$h=0$ を付け加えた点のテンソル平方が
元の六関係に入る。奇数位置の変更を戻すと標準曲線上の二つの非零切断の積が
恒等的に零となり、整性に反する。従って $T/(h)$ は可換完全交叉で各始点の次元が
$(1,3,3,1)$、最上次の積対合は完全である。

$h^3=0$ かつ $hr^2=hs^2=0$ なので、$(h)$ は
$h^e r^a s^b$ ($e=1,2$、$0\le a,b\le1$) で張られる。
$h$ を動かすときに $r,s$ の交換から生じる項も $h^3=0$ で消える。
よって $T$ の次数別上限は $(1,4,6,4,1)$、これまでの下限と一致し、各語が独立である。
特に $h^2rs\ne0$ である。$(h)$ の一次層、二次層の任意の非零元は
$h,r,s$ の相補語との積で最上次に届く。具体的には一次層ではまず $h$ を掛け、
残る $r,s$ の平方零の二文字で完全対合を取る。
$z$ は $(h)$ を消すが、$h^2=z\ell$ と書けば $z(\ell rs)=h^2rs\ne0$。
従って $T/(h)$ のsocleの持上げに対する $z$ の乗法は非零である。
$(h)$ の外の任意の斉次元も、まず完全交叉商の対合でsocleの持上げに移し、
$z$ を掛ければ最上次に届く。左右で同じ議論が使え、全相補積が完全となる。
これで有限Frobenius検査が通り、前節の三可換文字の命題で全次数AS正則性が従う。
\end{proof}

\begin{proof}[同型、次元と実在]
標準座標で二点の順序を保つ変換は $J$ を保ち、補文字をそれぞれ
$ax+\ell,by+m$ ($\ell,m\in J$) に移す。二点を通る割線を含む二次式の鉛筆から、
前出の跡零写像の二固有直線 $kr,ks$ が無順序で決まる。
交換する場合は表示した有限交換を先に施せる。各直線を保つときは
鉛筆の $x,y$ の係数比から $a=b$ である。$xy$ の式の変換が三式の包に入るには
具体的に
$m=m_r r+m_s s,\ \ell=-m_r r+m_s s$ となる。
これは $xm+y\ell$ の補文字を含む部分が $r(x-y),s(x+y)$ の包に入るためである。
従って $u,v$ の許される平行移動はそれぞれ $r,s$ の倍に限られ、
正規化 $f_R=g_S=0$ がこの自由度を消す。残る $h$ の持上げには $rh,sh$ の係数が生じるが、
少なくとも一つの $H$ 係数が非零なのでこれも零である。従って表示した倍率と交換だけが残る。
$h$ の符号反転の固定点は $h=0$ にある。滑らかな曲線の固定点で微分は $-1$、
他の座標比は不変なので $h$ はそこで単純に消える。固定点数は標準次数8、Hurwitz公式で商の種数は1である。

$u$ を保ち $v,s$ の符号を変えれば二点を交換して $\lambda$ の符号だけを変える。
従って述べた周期判定を得る。一般の台では
$f_S^2d_R/d_S^2,\ g_R^2d_S/d_R^2$ は倍率不変で、有限交換で互いの負に移る。
一般にはこの二数は互いの負でなく、交換を伴う同一視はない。
$d_R,d_S,f_S,g_R,d_H$ を保つ倍率は $a=b=e,\ c^2=e^2$ で $\lambda$ を変えない。
従って共通倍率を除いた一般安定化群は有限で、$\lambda\ne0$ の一般点は最小周期2である。

非空性は正規化前の $a=2$ の例
\[
 xy=r^2+2s^2+3h^2,\quad xr-2ry=2r^2+5s^2+h^2,\quad
 xs+2sy=3r^2-s^2+2h^2
\]
で確認できる。$y$ を再尺度化して平方完成と二平行移動を施せば定理の表示になる。
$h^2$ を消した二次式二本の射影三空間での交叉と、そこでの $h=0$ の分岐因子は滑らかである。
付随コードは標数101で両Jacobianの特異イデアルが全射影チャートで単位を含むことを確かめる。
次元定理で二交叉の次元はそれぞれ1,0であり、射影性によって標数零でも滑らかである。
この二重被覆が滑らかな完全交叉を与える。従って八次元係数空間の非空開集合が得られ、
一般軌道の三次元を引いて同型類の次元は5である。
\end{proof}
